% ===================================================================
%  QUADRO N.º 10 — O Mar. — O Porto.
%
%  Ceci n'est PAS une transcription : c'est une traduction, faite en
%  2026, en portugais d'aujourd'hui. Voir texto/pt/10-quadro-01.tex
%  pour la consigne, qui vaut pour les seize fichiers.
% ===================================================================

%%K t10-apar-1 apar
\VUsaut{4.0mm}
\VUtitre{15.0pt}{60}{QUADRO N.º 10}
\VUsaut{3.0mm}
\VUpk{11.4pt}{40}{O Mar. --- O Porto.}
\VUsaut{3.0mm}

%%K t10-01-1 p
1. --- No verão, enquanto uns vão buscar a calma e o fresco à montanha,
outros preferem ir à beira-mar. Foi o que fizemos no ano passado, porque
gostamos muito do \VUgras{oceano} \textsuperscript{(1)}.

%%K t10-02-1 p
2. --- Pela minha parte, sinto um grande prazer ao contemplar essa imensa
toalha de água que se estende até ao \VUgras{horizonte} \textsuperscript{(2)} e ao
ver partir para uma longa \VUgras{travessia} os enormes \VUgras{navios a
vapor} \textsuperscript{(3)} nos quais embarcam os viajantes que vão para a
América.

%%K t10-03-1 p
3. --- Por isso o meu pai, que conhece os meus gostos, escolhe sempre,
para passar as férias, uma praia muito tranquila situada perto de um dos
nossos grandes \VUgras{portos de guerra}.

%%K t10-03-2 p
Durante a nossa última estada, a \VUgras{esquadra} tinha vindo lançar
âncora atrás do enorme \VUgras{quebra-mar} \textsuperscript{(4)} que fecha e
protege o \VUgras{porto militar} \textsuperscript{(5)}, e pude admirar à minha
vontade os diferentes \VUgras{navios de guerra}, desde o pequeno
\VUgras{submarino} \textsuperscript{(10)}, que mergulha e desaparece sob as águas,
até ao enorme \VUgras{couraçado} \textsuperscript{(9)}, que até agora quase só
temia os ataques do \VUgras{torpedeiro} \textsuperscript{(8)}.

%%K t10-tit-2 sub
\VUsaut{3.0mm}
\VUpk{10.2pt}{40}{As embarcações pequenas.}
\VUsaut{3.0mm}

%%K t10-04-1 p
4. --- Este já sabia descobrir muito \VUgras{habilmente} o ponto fraco da
grossa couraça metálica para lhe cravar o perigoso \VUgras{torpedo}, cuja
explosão causa a perda do navio; mas era menos temível do que o
submarino, porque os contratorpedeiros lhe davam caça facilmente.

%%K t10-05-1 p
5. --- Pude também ver os rápidos \VUgras{cruzadores} \textsuperscript{(6)}
couraçados, quase tão invulneráveis como o verdadeiro couraçado, vários
\VUgras{transportes} \textsuperscript{(12)}, três ou quatro \VUgras{canhoneiras}
\textsuperscript{(7)} e, por fim, uma \VUgras{fragata} \textsuperscript{(11)}. \VUgras{O
espetáculo dessa frota, quando manobra para levantar âncoras, é dos mais
interessantes.}

%%K t10-06-1 p
6. --- Que diferença de construção entre essas enormes massas de aço e os
elegantes \VUgras{brigues de três mastros} ou os finos \VUgras{veleiros}
\textsuperscript{(13)}, que ainda se usam na marinha mercante, e que eu via
entrar no porto a todo o pano quando ia passear pelo \VUgras{molhe}
\textsuperscript{(14)}! É verdade que estes perdiam muito das suas vantagens
quando o \VUgras{vento} era contrário. Tinham de recolher todas as velas
para entrar na doca, e era preciso um \VUgras{rebocador}
\textsuperscript{(16)} que os fosse buscar para os levar ao cais como simples
\VUgras{barcaças} \textsuperscript{(17)}.

%%K t10-tit-3 sub
\VUsaut{3.0mm}
\VUpk{10.2pt}{40}{O porto comercial.}
\VUsaut{3.0mm}

%%K t10-07-1 p
7. --- O interessante do porto de que falo é que não compreende só um
porto de guerra, criado artificialmente com obras gigantescas, mas ainda
um maravilhoso \VUgras{porto de comércio} situado na \VUgras{foz} de um
grande rio.

%%K t10-08-1 p
8. --- A língua de terra que separa as duas docas está fortificada e
defendida por uma série de \VUgras{baterias} e \VUgras{baluartes} que
entram pelo mar dentro. É preciso uma licença especial para passear pelas
\VUgras{muralhas} \textsuperscript{(15)}. Eu tinha-a obtido sem dificuldade por
intermédio do meu tio, que é engenheiro dos \VUgras{estaleiros}
\textsuperscript{(18)}, e fui visitar os navios que se constroem sob vastos
alpendres.

%%K t10-09-1 p
9. --- Até assisti ao lançamento à água de um couraçado e lembro-me do
momento de emoção que sentiu toda a multidão quando a enorme massa,
abandonada a si mesma, começou a deslizar pela sua carreira e veio
flutuar sobre a água do rio.

%%K t10-10-1 p
10. --- Para ir aos estaleiros atravessa-se o \VUgras{campo de manobras}
\textsuperscript{(19)}, onde as tropas da guarnição vêm cada dia fazer o
exercício, e passa-se por uma \VUgras{passarela} \textsuperscript{(20)} de ferro
que une a esplanada ao \VUgras{arsenal} \textsuperscript{(21)}.

%%K t10-tit-4 sub
\VUsaut{3.0mm}
\VUpk{10.2pt}{40}{O Arsenal e as docas secas.}
\VUsaut{3.0mm}

%%K t10-11-1 p
11. --- Com o meu tio visitei esse grande estabelecimento cheio de
\VUgras{canhões} \textsuperscript{(22)}, \VUgras{balas} \textsuperscript{(23)} e
\VUgras{obuses}. Vi também a \VUgras{fábrica metalúrgica} \textsuperscript{(24)} na
qual se fabricam as \VUgras{caldeiras} \textsuperscript{(25)} dos navios a vapor.

%%K t10-12-1 p
12. --- Muito perto estão as \VUgras{docas secas} \textsuperscript{(26)} --- docas
de carena --- e as curiosíssimas \VUgras{docas flutuantes}
\textsuperscript{(27)}, nas quais pude ver de perto, num navio em reparação, a
\VUgras{hélice} \textsuperscript{(28)}, a \VUgras{quilha} \textsuperscript{(29)} e o
\VUgras{leme} \textsuperscript{(30)} de uma nau.

%%K t10-13-1 p
13. --- O porto de comércio é tão interessante de estudar como o porto de
guerra, e passei muitas horas a vaguear pelos cais onde estão amarrados
os \VUgras{vapores de rodas} \textsuperscript{(31)} que sobem o rio, e a ver
funcionar a \VUgras{cabrea} \textsuperscript{(32)}, que levanta como plumas cargas
enormes.

%%K t10-tit-5 sub
\VUsaut{4.0mm}
\VUpk{10.2pt}{40}{O Piloto.}
\VUsaut{3.0mm}

%%K t10-14-1 p
14. --- Admirava os \VUgras{transatlânticos} \textsuperscript{(34)} que saíam da
enseada com as suas \VUgras{bandeiras} \textsuperscript{(35)} ao vento, conduzidos
por um hábil \VUgras{piloto} \textsuperscript{(36)} que sabia maravilhosamente
seguir o \VUgras{canal} marcado por \VUgras{boias} \textsuperscript{(40)} e
\VUgras{balizas}, evitando as \VUgras{barcaças} \textsuperscript{(41)} que se
governam com tanta dificuldade com a sua única vela quadrada. Distinguia
na \VUgras{proa} \textsuperscript{(37)} os \VUgras{camarotes} \textsuperscript{(39)} de
segunda classe e, na \VUgras{popa} \textsuperscript{(38)}, os salões dos
passageiros de primeira.

%%K t10-15-1 p
15. --- Às vezes assistia também à carga de um \VUgras{navio de carga}
\textsuperscript{(42)}, no qual se amontoavam as mercadorias do \VUgras{armazém}
\textsuperscript{(43)} e da \VUgras{estação marítima} \textsuperscript{(44)}, trazidas
pelos numerosos comboios da \VUgras{linha dos cais} \textsuperscript{(45)}.

%%K t10-tit-6 sub
\VUsaut{3.0mm}
\VUpk{10.2pt}{40}{Remar por gosto.}
\VUsaut{3.0mm}

%%K t10-16-1 p
16. --- Muitas vezes remava na \VUgras{doca de flutuação} \textsuperscript{(53)},
onde se refugiam as pequenas \VUgras{embarcações} \textsuperscript{(46)}, e era
uma alegria para mim manejar o \VUgras{remo} \textsuperscript{(47)}. Subia ao
\VUgras{convés} \textsuperscript{(48)} de um \VUgras{barco à vela}
\textsuperscript{(49)}, trepava com o auxílio do \VUgras{cordame}
\textsuperscript{(52)} pelo \VUgras{mastro} \textsuperscript{(51)} até às \VUgras{vergas}
\textsuperscript{(50)}, e depois retomava o meu passeio de \VUgras{iole}
\textsuperscript{(54)} entre os pesados \VUgras{barcos de pesca}
\textsuperscript{(55)}, onde secam as \VUgras{redes} \textsuperscript{(56)}.

%%K t10-17-1 p
17. --- Pelo \VUgras{cais} \textsuperscript{(58)} desfilavam diante de mim as
\VUgras{mercadorias} \textsuperscript{(59)} mais diversas, metidas em
\VUgras{caixas} \textsuperscript{(60)} ou em \VUgras{fardos} \textsuperscript{(61)}.
Formavam o \VUgras{carregamento} \textsuperscript{(63)} das barcaças, e potentes
\VUgras{guindastes} \textsuperscript{(62)} levantavam-nas e deixavam-nas em
\VUgras{camiões} \textsuperscript{(64)} enquanto os \VUgras{transportadores} faziam
a declaração e pagavam os direitos na \VUgras{alfândega} \textsuperscript{(65)}.

%%K t10-18-1 p
18. --- Por fim, quando chegava à \VUgras{ponte giratória}
\textsuperscript{(66)} ou à \VUgras{comporta} \textsuperscript{(69)}, passando diante da
\VUgras{draga} \textsuperscript{(68)} que limpa o \VUgras{canal} \textsuperscript{(67)},
desembarcava no molhe, junto ao \VUgras{semáforo} \textsuperscript{(70)}.

%%K t10-19-1 p
19. --- É do outro lado desse molhe que atracam as \VUgras{chalupas}
\textsuperscript{(71)} que trazem a terra os oficiais da esquadra. É um belo
espetáculo ver os marinheiros levantar os seus remos a compasso enquanto
a embarcação, levada pela \VUgras{onda} \textsuperscript{(72)}, aborda sem esforço
a escada do desembarcadouro. É aqui que melhor se pode observar a
\VUgras{maré} e o movimento do \VUgras{fluxo} e do \VUgras{refluxo} na
\VUgras{praia} \textsuperscript{(73)}.

%%K t10-tit-7 sub
\VUsaut{3.0mm}
\VUpk{10.2pt}{40}{A Messe dos oficiais.}
\VUsaut{3.0mm}

%%K t10-20-1 p
20. --- Para voltar a casa tomava o \VUgras{elétrico} \textsuperscript{(74)}, cuja
\VUgras{paragem} \textsuperscript{(75)} está no \VUgras{poste} que há diante da
messe dos oficiais.

%%K t10-20-2 p
Da \VUgras{varanda} \textsuperscript{(76)} da messe, enfeitada com \VUgras{âncoras}
\textsuperscript{(77)}, canhões e balas, via muitas vezes uma brilhante reunião
de oficiais de marinha: \VUgras{primeiros-tenentes} \textsuperscript{(78)} com a
sua espada, \VUgras{capitães de artilharia} \textsuperscript{(79)} e de
\VUgras{infantaria} \textsuperscript{(80)} de marinha com o seu \VUgras{sabre}.
Atrás deles estava um \VUgras{marinheiro} \textsuperscript{(81)} que lhes trazia
um \VUgras{óculo} quando queriam distinguir o que se passava ao longe.

%%K t10-21-1 p
21. --- Via também jovens \VUgras{segundos-tenentes} \textsuperscript{(82)} que
recebiam as ordens do seu \VUgras{comandante} \textsuperscript{(83)} ou do
\VUgras{almirante} \textsuperscript{(84)}, enquanto um \VUgras{cabo}
\textsuperscript{(85)} esperava para as levar a bordo.

%%K t10-22-1 p
22. --- Dessa varanda a vista é magnífica: domina-se toda a praia com as
suas \VUgras{tendas} \textsuperscript{(86)} e as suas \VUgras{barracas}
\textsuperscript{(87)}, e vê-se, à hora do banho, os \VUgras{banhistas}
\textsuperscript{(88)} que folgam alegremente diante do \VUgras{casino}
\textsuperscript{(89)}.

%%K t10-23-1 p
23. --- Na ponta da praia, a costa forma uma pequena \VUgras{península}
\textsuperscript{(91)} rochosa, onde vem quebrar a \VUgras{rebentação}
\textsuperscript{(92)} e sobre a qual está construído um \VUgras{farol}
\textsuperscript{(93)} que indica de noite o caminho aos navegantes. Essa
península só está unida a terra por um \VUgras{istmo} \textsuperscript{(90)} muito
estreito.

%%K t10-tit-8 sub
\VUsaut{3.0mm}
\VUpk{10.2pt}{40}{Vista geral da costa.}
\VUsaut{3.0mm}

%%K t10-24-1 p
24. --- A \VUgras{arriba} \textsuperscript{(94)} prolonga-se, recortada pelo mar,
que nela desenha caprichosamente uma série de \VUgras{baías}
\textsuperscript{(95)} e \VUgras{promontórios} (cabos) que a tornam dos mais
pitorescos.

%%K t10-24-2 p
No \VUgras{planalto} \textsuperscript{(96)}, que domina o \VUgras{estreito}
\textsuperscript{(98)}, há um \VUgras{forte} \textsuperscript{(97)} muito importante e
alguns chalés.

%%K t10-25-1 p
25. --- Esse estreito separa do continente uma \VUgras{ilha}
\textsuperscript{(99)} muito perigosa, rodeada de \VUgras{escolhos}
\textsuperscript{(100)} e \VUgras{recifes}, onde barcos e até grandes navios
encalham às vezes. Apesar da prontidão dos socorros e da rápida chegada
dos \VUgras{botes salva-vidas}, esses \VUgras{naufrágios} são terríveis e,
nas fortes \VUgras{tempestades} de equinócio, vários navios se perderam
com homens e mercadorias.

%%K t10-25-2 p
Apesar dessa vizinhança, \VUgras{o porto} é muito frequentado pelos
marinheiros.
\VUsaut{6.0mm}
\VUfilet{22mm}
